% !TeX encoding = UTF-8
% Artículo para la memoria del Congreso Internacional de Humanidades,
% Ingeniería y Tecnología (CIHITEC 2026). Versión condensada de la tesis
% original (docs/tesisss.tex), adaptada a las Normas de publicación:
% Times New Roman 12 pt, interlineado 1.5, márgenes 2.5 cm, APA 7, 8-15 págs.
\documentclass[12pt,letterpaper]{article}
\usepackage[utf8]{inputenc}
\usepackage[T1]{fontenc}
\usepackage{mathptmx} % Times New Roman en texto y matemáticas
\usepackage[spanish,mexico]{babel}
\usepackage{graphicx}
\usepackage{booktabs}
\usepackage{amsmath}
\usepackage{amssymb}
\usepackage{placeins}   % \FloatBarrier
\usepackage[margin=2.5cm]{geometry}
\usepackage[hidelinks]{hyperref}
\usepackage[natbibapa]{apacite} % Citas y referencias en estilo APA (autor-año)

\renewcommand{\baselinestretch}{1.5} % Interlineado 1.5
\setcounter{secnumdepth}{3}          % Numeración 1 / 1.1 / 1.1.1
\graphicspath{{img/}{./}}

\title{\textbf{Logística de última milla en redes viales reales: un enfoque
reproducible para el MDRP en La Paz, B.C.S.}}
\author{Alan José Marcos Sánchez Martínez\\
	\small Instituto Tecnológico de La Paz, Tecnológico Nacional de México (TecNM)\\
	\small alansama2906@gmail.com}
\date{}

\begin{document}
\maketitle

\begin{abstract}
\noindent La logística de última milla concentra entre el 30\,\% y el 60\,\% del
costo de una operación de entrega y, en el segmento de comida preparada, se
agrava por la urgencia y la demanda dinámica. La mayoría de los estudios sobre el
\textit{Meal Delivery Routing Problem} (MDRP) evalúan sus políticas de despacho
con distancias euclidianas, que subestiman los trayectos reales entre un 17\,\% y
un 40\,\% y comprometen la validez de los tiempos de entrega prometidos. Este
trabajo evalúa la eficiencia operativa de la política anticipatoria de horizonte
rodante (\textit{Rolling Horizon}, RH) frente a una política reactiva
\textit{First-Come, First-Served} (FCFS) para el MDRP sobre la red vial real de
La Paz, Baja California Sur. Se desarrolló un entorno de simulación de eventos
discretos reproducible, basado exclusivamente en software libre (OSRM,
OpenStreetMap y Docker), que sustituye las distancias euclidianas por tiempos de
viaje reales. La demanda se generó mediante un proceso de Poisson no homogéneo
bimodal calibrado con el conjunto de datos LaDe (aproximadamente 1\,038 pedidos y
25 restaurantes) y ambas políticas se compararon en cuatro escenarios de
saturación de flota con siete réplicas válidas, contrastadas con la prueba de
Mann-Whitney~$U$. Los resultados revelan dos regímenes: con flota holgada FCFS es
más rápida, pero bajo saturación media y alta RH reduce el \textit{Click-to-Door}
promedio hasta un 57.2\,\% y su percentil~90 hasta un 77.1\,\%, además de eliminar
los pedidos no entregados. Se identifica un punto de cruce entre 26 y 35 pedidos
por repartidor a partir del cual la anticipación supera a la asignación reactiva.

\vspace{0.4cm}
\noindent\textbf{Palabras clave:} MDRP; horizonte rodante; ruteo dinámico de
vehículos; logística de última milla; OSRM; simulación de eventos discretos;
software libre.
\end{abstract}

% ======================================================================
\section{Introducción}
% ======================================================================

La entrega de última milla ---el segmento final desde el punto de origen hasta el
domicilio del cliente--- es el componente más costoso y complejo de la logística
urbana: puede representar entre el 30\,\% y el 60\,\% del costo total de un envío
\citep{Jazemi2023}. En el caso de la comida preparada, este segmento se vuelve
particularmente exigente por la urgencia (ventanas de 30 a 60 minutos), la
perecibilidad del producto y una demanda dinámica e impredecible que obliga a
reoptimizar continuamente las asignaciones.

El problema de enrutamiento de vehículos (VRP), formulado originalmente por
\citet{Dantzig1959} como generalización del problema del agente viajero, ha
evolucionado durante cinco décadas para incorporar restricciones de capacidad,
ventanas de tiempo y múltiples depósitos \citep{Laporte2009VRPSurvey}. Cuando la
demanda se revela en tiempo real, el problema se torna dinámico
\citep{Larsen2000DVRP,Pillac2013}. En esta familia, \citet{Reyes2018}
introdujeron el \textit{Meal Delivery Routing Problem} (MDRP), que combina
múltiples puntos de recolección, tiempos de preparación, alta urgencia y
agrupamiento dinámico de pedidos (\textit{bundling}).

Un problema transversal a esta literatura es la manera de estimar los tiempos de
viaje. Las plataformas comerciales (Google Maps, Mapbox) ofrecen precisión pero
imponen costos de licenciamiento y operan como caja negra, comprometiendo la
reproducibilidad científica \citep{Pereira2021}. La alternativa académica
tradicional ---la distancia euclidiana--- es gratuita y reproducible, pero
\citet{Boyaci2021} demostraron, sobre 96 instancias en 12 ciudades, que las rutas
reales son entre un 17\,\% y un 40\,\% más largas que la línea recta, lo que
invalida las estimaciones de tiempo de arribo y puede producir rutas
operativamente infactibles.

Esta brecha se agudiza en ciudades medianas de topología irregular como La Paz,
Baja California Sur, donde la ensenada obliga a rodeos extensos y la estructura
vial depende de pocas arterias estructurantes. Además, \citet{Reyes2018}
validaron sus políticas de horizonte rodante en instancias sintéticas de grandes
urbes con alta conectividad y distancias euclidianas; se desconoce el grado de
degradación de estas heurísticas al trasladarlas a un entorno de baja
conectividad y alta penalización por error.

De lo anterior surge la pregunta de investigación: \textit{¿en qué medida una
política de despacho anticipatoria supera a una reactiva cuando el MDRP se evalúa
sobre la red vial real de una ciudad mediana de topología hostil?} El
\textbf{objetivo general} es evaluar la eficiencia operativa de la política
\textit{Rolling Horizon} frente a la estrategia reactiva FCFS mediante un entorno
de simulación reproducible que integre las restricciones topológicas reales de La
Paz y patrones de demanda estocásticos. La \textbf{hipótesis} sostiene que la
política anticipatoria, alimentada con tiempos de viaje reales, superará
significativamente a la reactiva en eficiencia operativa y nivel de servicio,
usando exclusivamente software de código abierto.

El resto del documento se organiza como sigue. La Sección~2 presenta el marco
teórico (formulación del MDRP, heurística de horizonte rodante, enrutamiento
sobre grafos viales y generación estocástica de demanda). La Sección~3 describe la
metodología y el diseño experimental. La Sección~4 presenta y discute los
resultados. La Sección~5 concluye y plantea líneas de trabajo futuro.

% ======================================================================
\section{Marco teórico}
% ======================================================================

\subsection{Formulación del MDRP}

El MDRP \citep{Reyes2018} se define sobre tres conjuntos: restaurantes
$\mathcal{R}$ (con ubicación fija), órdenes $\mathcal{O}$ y repartidores
(\textit{couriers}) $\mathcal{C}$. Cada orden $o$ se caracteriza por su
restaurante de origen $r_o$, su instante de colocación $a_o$, su instante de
disponibilidad o \textit{ready time} $e_o \ge a_o$ y su ubicación de entrega
$\ell_o$; la información de la orden se revela al sistema únicamente en $a_o$
(llegada dinámica). Toda la información de restaurantes y couriers se conoce a
priori. El problema consiste en determinar rutas factibles que completen la
recolección y entrega de las órdenes optimizando una o varias métricas de
desempeño.

Dos suposiciones estructurales son centrales. Primero, el \textbf{agrupamiento}
(\textit{bundling}): órdenes de un mismo restaurante pueden combinarse en una ruta
con múltiples entregas, cuyo \textit{ready time} es el máximo de sus órdenes
constituyentes. Segundo, la \textbf{restricción de \textit{ready time}}: un
courier no puede recoger una orden antes de que la comida esté lista. La
remuneración sigue un esquema de compensación mínima garantizada:
\begin{equation}
	\text{Pago}(c) = \max\!\left(p_1\, n_c,\; p_2\,(l_c - e_c)\right),
	\label{eq:comp}
\end{equation}
donde $n_c$ es el número de órdenes entregadas por el courier $c$, $l_c-e_c$ la
duración de su turno, $p_1$ el pago por orden y $p_2$ el pago mínimo por hora.

Las métricas primarias relevantes son el \textbf{\textit{Click-to-Door}} (CtD,
tiempo entre colocación y entrega, $\text{CtD}_o = t_{delivery,o}-a_o$), el
\textbf{\textit{Ready-to-Pickup}} (RtP, espera de la comida lista) y el
\textbf{\textit{Ready-to-Door}} (RtD), además de la utilización del courier y el
costo por orden. El MDRP pertenece a los problemas dinámicos de recolección y
entrega \citep{Pillac2013} y se distingue por combinar múltiples puntos de
recolección, alta urgencia ($\tau\approx 40$~min) y dinamismo total, lo que
descarta los métodos exactos y favorece heurísticas basadas en horizontes
rodantes.

\subsection{Heurística de horizonte rodante}

El algoritmo de \citet{Reyes2018} ejecuta, cada $f$~minutos, un ciclo de
agrupamiento, asignación y compromiso sobre las órdenes cuyo \textit{ready time}
cae dentro de un horizonte de asignación $\Delta_u$. Para equilibrar la
utilización de la flota con la frescura de las entregas, el \textbf{tamaño
objetivo de agrupamiento} al instante $t$ se define como la razón entre las
órdenes próximas a estar listas y los couriers próximos a estar disponibles:
\begin{equation}
	Z_t = \frac{|\{o \in \mathcal{O}_t : e_o \le t + \Delta_1\}|}
	           {|\{c \in \mathcal{D}_t : e^{avail}_c \le t + \Delta_2\}|},
	\qquad \Delta_1,\Delta_2 > 0,
	\label{eq:zt}
\end{equation}
donde $\mathcal{D}_t$ es el conjunto de couriers en turno. Cuando $Z_t<1$ hay más
couriers que órdenes y el comportamiento converge a asignaciones individuales;
cuando $Z_t\gg 1$ el sistema está saturado y se favorecen agrupamientos mayores.

Determinado $Z_t$, se construyen agrupamientos por restaurante mediante inserción
paralela con mejora local (\textit{remove-reinsert}), y se clasifican en tres
grupos de prioridad según su urgencia respecto al objetivo CtD $\tau$. La
asignación se resuelve secuencialmente por grupo como un \textbf{problema de
asignación lineal}: para cada par agrupamiento-courier $(s,d)$ se define el peso
\begin{equation}
	w_{s,d} = \frac{N_s}{\max_{o \in s}\{\delta^s_{o,d}\} - e_d}
	        - \theta\!\left(\pi_{s,d} - \max_{o \in s}\{e_o\}\right),
	\label{eq:matching}
\end{equation}
donde $N_s$ es el número de órdenes del agrupamiento, $\delta^s_{o,d}$ el tiempo
de entrega de la orden $o$, $e_d$ el momento en que el courier $d$ queda
disponible, $\pi_{s,d}$ el tiempo de recolección y $\theta>0$ penaliza el retraso
respecto al \textit{ready time}. El primer término premia la eficiencia (órdenes
por unidad de tiempo) y el segundo la frescura. El problema se resuelve en
$O(n^3)$ con el algoritmo húngaro \citep{Kuhn1955}.

Finalmente, una estrategia de \textbf{compromiso aditivo en dos etapas}
\citep{Reyes2018} distingue entre compromiso final (el courier recoge y entrega
en la época actual) y compromiso parcial o \textit{inbound} (el courier se
desplaza hacia el restaurante y su agrupamiento se confirma en la siguiente
época). El compromiso parcial actúa como preposicionamiento y conserva
flexibilidad ante nueva información.

\subsection{Enrutamiento sobre grafos viales}

Para eliminar el error de la aproximación euclidiana, este trabajo sustituye las
distancias en línea recta por consultas sobre el grafo vial real mediante la
cadena OpenStreetMap~$\rightarrow$~OSRM. OpenStreetMap (OSM) modela la
infraestructura mediante nodos, vías y relaciones \citep{Haklay2010,Boeing2021},
de las que se deriva un grafo $G=(V,E,c)$ con costos definidos por la distancia y
la velocidad permitida de cada segmento.

El cálculo de caminos mínimos se acelera con \textit{Contraction Hierarchies} (CH)
\citep{Geisberger2008}, técnica de preprocesamiento que contrae los nodos por
orden de importancia e inserta aristas de atajo, permitiendo resolver consultas
mediante búsqueda bidireccional en milisegundos. \textbf{OSRM} (\textit{Open
Source Routing Machine}) implementa CH sobre datos OSM en tres etapas
(\texttt{osrm-extract}, \texttt{osrm-contract} y \texttt{osrm-routed}), con
tiempos de respuesta típicos inferiores a 5~ms para grafos regionales. La
sustitución de la distancia euclidiana por tiempos OSRM tiene tres consecuencias:
los tiempos dejan de ser simétricos ($T(a,b)\neq T(b,a)$) por los sentidos de
circulación, la desigualdad triangular se cumple de forma natural sobre la red y
se elimina el factor de conversión euclidiano de \citet{Reyes2018}.

\subsection{Generación estocástica de demanda}

Las llegadas de pedidos se modelan mediante un \textbf{proceso de Poisson no
homogéneo} \citep{Ross2014}, en el que el número de eventos en un intervalo sigue
una distribución de Poisson cuyo parámetro es la integral de una función de
intensidad $\lambda(t)$. Los servicios de comida exhiben dos picos diarios
(almuerzo y cena) \citep{Wu2023LaDe}, por lo que la intensidad se representa como
la superposición de una base constante y dos componentes gaussianas:
\begin{equation}
	\lambda(t) = \lambda_0
	+ A_m \exp\!\left(-\frac{(t-\mu_m)^2}{2\sigma^2}\right)
	+ A_e \exp\!\left(-\frac{(t-\mu_e)^2}{2\sigma^2}\right),
	\label{eq:intensity}
\end{equation}
con base $\lambda_0$, centros $\mu_m,\mu_e$, amplitudes $A_m,A_e$ y dispersión
$\sigma$. La propiedad de discretización de Poisson permite generar instancias
mediante un simple muestreo minuto a minuto. Las métricas de desempeño se agregan
con la media y los percentiles 50, 90 y 95; el percentil 90 es especialmente
relevante porque captura la experiencia de los clientes en la cola de la
distribución.

% ======================================================================
\section{Metodología}
% ======================================================================

\subsection{Arquitectura del sistema}

El entorno de simulación se organiza en tres capas funcionales desacopladas:
(1)~una capa de \textbf{enrutamiento} basada en OSRM~v5.27.1 desplegado en un
contenedor Docker \citep{Merkel2014} que opera sobre \texttt{localhost} con el
algoritmo CH sobre el extracto OSM de La Paz; (2)~una capa de \textbf{simulación}
de eventos discretos implementada en Python, que consulta tiempos y rutas al
servidor OSRM vía HTTP; y (3)~una capa de \textbf{análisis} que calcula los
indicadores y genera las tablas y figuras. La contenerización garantiza que el
entorno de cálculo sea idéntico en cualquier máquina anfitriona, y la operación
\textit{offline} elimina el consumo de APIs comerciales y las barreras económicas
para la replicación.

\subsection{Generación de instancias sintéticas}

Cada instancia contiene aproximadamente 1\,038 pedidos distribuidos en una ventana
operativa diaria y asociados a 25 restaurantes. Los instantes de colocación se
muestrean del proceso de Poisson no homogéneo bimodal de la
Ecuación~\eqref{eq:intensity}, calibrado con los picos de demanda del conjunto de
datos LaDe \citep{Wu2023LaDe}. Las ubicaciones de restaurantes y clientes se
muestrean dentro del polígono urbano de La Paz, validando que cada punto obtenga
respuesta del servicio de enrutamiento. Los tiempos de preparación se modelan como
$T_{prep}\sim\max(5,\mathcal{N}(12,3^2))$~min. El tamaño de flota se dimensiona
automáticamente según
\begin{equation}
	|\mathcal{C}| = \max\!\left(n_{\min},\; \lceil |\mathcal{O}|/r \rceil\right),
	\label{eq:fleet}
\end{equation}
donde $r$ es el número objetivo de pedidos por repartidor y $n_{\min}=10$. El uso
de semillas aleatorias garantiza la reproducibilidad exacta de las instancias.

\subsection{Motor de simulación y políticas de despacho}

La simulación avanza en \textbf{épocas de reoptimización} de $f=5$~min. En cada
época, cada orden transita por una máquina de estados (colocada
$\rightarrow$ lista $\rightarrow$ asignada $\rightarrow$ recogida
$\rightarrow$ entregada) y los tiempos de viaje se obtienen de OSRM. Se comparan
dos políticas:

\begin{itemize}
	\item \textbf{FCFS (línea base):} asignación reactiva e inmediata de cada
	      pedido al courier disponible más cercano al restaurante, sin
	      consolidación de órdenes.
	\item \textbf{Rolling Horizon (RH):} implementación de la heurística de la
	      Sección~2.2, con cálculo del tamaño objetivo de agrupamiento
	      (Ecuación~\ref{eq:zt}), generación de agrupamientos por restaurante,
	      clasificación por prioridad, asignación lineal
	      (Ecuación~\ref{eq:matching}) resuelta con el algoritmo húngaro
	      \citep{Kuhn1955} y compromiso en dos etapas.
\end{itemize}

Los parámetros del algoritmo se mantienen constantes: frecuencia de optimización
$f=5$~min, horizonte $\Delta_u=20$~min, objetivo SLA $\tau=40$~min, tiempo de
servicio $s=4$~min y tamaño máximo de agrupamiento $B_{max}=4$. La compensación
usa $p_1=\$10$ por pedido y un piso de $p_2=\$15$ por hora.

\subsection{Diseño experimental}

La variable de tratamiento es la \textbf{razón de saturación} $r$ (pedidos por
repartidor), que dimensiona la flota mediante la Ecuación~\eqref{eq:fleet}
manteniendo fijo el volumen de pedidos. Se definen cuatro escenarios que cubren
desde operación holgada hasta saturación alta (Tabla~\ref{tab:escenarios}).

\begin{table}[htbp]
\centering
\caption{Escenarios experimentales de saturación de flota ($|\mathcal{O}|\approx
1\,038$ pedidos). Fuente: elaboración propia.}
\label{tab:escenarios}
\begin{tabular}{@{}clccl@{}}
\toprule
Escenario & Nivel & $r$ & $|\mathcal{C}|$ & Descripción \\
\midrule
S1 & Baja     & 18 & 58 & Línea base; capacidad holgada \\
S2 & Moderada & 26 & 40 & Reducción moderada de flota \\
S3 & Media    & 35 & 30 & Punto medio de saturación \\
S4 & Alta     & 52 & 20 & Carga elevada por repartidor \\
\bottomrule
\end{tabular}
\end{table}

Para cada combinación de escenario y semilla, ambas políticas se ejecutan sobre la
\emph{misma} instancia de pedidos y flota (comparación \textbf{pareada}), lo que
aísla el efecto de la política de asignación. Se emplearon diez semillas
($s\in\{2025,\ldots,2034\}$), de las cuales siete produjeron resultados válidos;
las tres restantes se descartaron porque la totalidad de sus órdenes permaneció en
estado \emph{ready} sin asignación, atribuible a una interrupción del servicio
OSRM. Los valores reportados corresponden al promedio de las siete réplicas
válidas, generando $4\times 7\times 2 = 56$ corridas.

El desempeño se cuantifica mediante 21 indicadores agrupados en cinco dimensiones:
calidad de servicio (CtD, RtP, RtD, cumplimiento SLA y exceso sobre el objetivo o
\textit{CTD Overage}), cobertura (\% de pedidos no entregados), eficiencia
operativa (pedidos por repartidor por hora, utilización y distancia por pedido),
eficiencia de agrupamiento (tamaño promedio de agrupamiento) y costo (costo por
pedido, derivado de la Ecuación~\ref{eq:comp}). Para validar la significancia de
las diferencias se aplica la prueba no paramétrica de Mann-Whitney~$U$
\citep{MannWhitney1947} con corrección de Benjamini-Hochberg y correlación
rank-biserial $r$ como tamaño del efecto \citep{Cohen1988}, complementada con la
prueba ómnibus de Kruskal-Wallis. Se eligió Mann-Whitney~$U$ porque la prueba de
Shapiro-Wilk rechazó la normalidad de los KPIs y por su robustez ante valores
atípicos.

% ======================================================================
\section{Resultados y discusión}
% ======================================================================

\subsection{Calidad de servicio}

La Tabla~\ref{tab:res_ctd} presenta los indicadores de calidad de servicio
promediados sobre las siete réplicas válidas. El \textit{Click-to-Door} (CtD) es
la métrica principal de calidad percibida por el cliente.

\begin{table}[htbp]
\centering
\small
\caption{Indicadores de calidad de servicio: promedio de 7 réplicas. Fuente:
elaboración propia.}
\label{tab:res_ctd}
\begin{tabular}{@{}l rr r rr r rr r@{}}
\toprule
& \multicolumn{3}{c}{Avg CtD (min)} & \multicolumn{3}{c}{P90 CtD (min)} & \multicolumn{3}{c}{SLA (\%)} \\
\cmidrule(lr){2-4} \cmidrule(lr){5-7} \cmidrule(lr){8-10}
Esc. & FCFS & RH & $\Delta$\% & FCFS & RH & $\Delta$\% & FCFS & RH & $\Delta$\% \\
\midrule
S1 & 10.10 & 20.38 & $+101.8$ & 13.10  & 31.51 & $+140.5$ & 100.0 & 99.2 & $-0.8$  \\
S2 & 16.68 & 22.74 & $+36.3$  & 32.73  & 35.46 & $+8.3$   & 91.7  & 97.0 & $+5.8$  \\
S3 & 35.79 & 25.52 & $-28.7$  & 117.19 & 39.94 & $-65.9$  & 76.3  & 90.2 & $+18.2$ \\
S4 & 75.41 & 32.26 & $-57.2$  & 228.01 & 52.27 & $-77.1$  & 61.2  & 78.3 & $+27.9$ \\
\bottomrule
\end{tabular}
\end{table}

En el escenario de baja saturación (S1), con oferta holgada de repartidores, FCFS
logra tiempos menores que RH (10.10 vs 20.38~min), porque despacha cada pedido de
inmediato al courier más cercano mientras que RH introduce la latencia del ciclo
de optimización y la espera para formar agrupamientos. No obstante, ambas
políticas se mantienen muy por debajo del objetivo SLA de 40~min. En S2 la brecha
se reduce ($+36.3\,\%$) y, notablemente, el cumplimiento SLA se invierte a favor de
RH (97.0\,\% vs 91.7\,\%): antes incluso de superar a FCFS en promedio, la
coordinación de agrupamientos ya reduce los pedidos que exceden el umbral.

A partir de S3, RH supera decisivamente a FCFS: el CtD promedio baja un 28.7\,\%
y el percentil~90 cae de 117.19 a 39.94~min ($-65.9\,\%$), comprimiendo la cola
derecha de la distribución. En S4 la ventaja es máxima: el CtD promedio se reduce
un 57.2\,\% y el P90 un 77.1\,\%; además, bajo FCFS un 11.5\,\% de los pedidos
queda sin entregar al cierre, mientras que bajo RH se entrega la totalidad
(0\,\%). La Figura~\ref{fig:ctd_boxplot} ilustra el colapso progresivo de FCFS en
S3--S4 (cajas con alta dispersión y valores atípicos por encima de 300~min) frente
a la distribución compacta de RH.

\begin{figure}[htbp]
\centering
\includegraphics[width=0.9\textwidth]{img/fig_ctd_boxplot.png}
\caption{Distribución del \textit{Click-to-Door} por escenario y política. La
línea punteada indica el objetivo SLA de $\tau=40$~min. Fuente: elaboración
propia.}
\label{fig:ctd_boxplot}
\end{figure}

\FloatBarrier

\subsection{Eficiencia operativa y agrupamiento}

La Tabla~\ref{tab:res_eficiencia} resume la eficiencia operativa. El
\textit{throughput} (pedidos por repartidor por hora) bajo RH es ligeramente
superior en S1--S3 y se amplía a $+17.7\,\%$ en S4, gracias al agrupamiento: al
asignar varios pedidos por viaje, cada turno produce más entregas. La utilización
del courier, en cambio, es \emph{menor} bajo RH en S2--S4; esto no indica
ociosidad sino mayor eficiencia kilométrica, pues los agrupamientos reducen la
distancia por pedido hasta un 39.9\,\% en S4 (de 5.41 a 3.25~km), con el
consiguiente ahorro de combustible y tiempo de conducción.

\begin{table}[htbp]
\centering
\small
\caption{Indicadores de eficiencia operativa: promedio de 7 réplicas. Fuente:
elaboración propia.}
\label{tab:res_eficiencia}
\begin{tabular}{@{}l rr r rr r rr r@{}}
\toprule
& \multicolumn{3}{c}{Ped./Rep./Hr} & \multicolumn{3}{c}{Utilización (\%)} & \multicolumn{3}{c}{Dist./Pedido (km)} \\
\cmidrule(lr){2-4} \cmidrule(lr){5-7} \cmidrule(lr){8-10}
Esc. & FCFS & RH & $\Delta$\% & FCFS & RH & $\Delta$\% & FCFS & RH & $\Delta$\% \\
\midrule
S1 & 1.43 & 1.50 & $+4.9$  & 18.0 & 20.9 & $+16.1$ & 4.49 & 4.73 & $+5.3$  \\
S2 & 2.07 & 2.14 & $+3.4$  & 27.3 & 26.9 & $-1.5$  & 4.77 & 4.23 & $-11.3$ \\
S3 & 2.77 & 2.90 & $+4.7$  & 38.4 & 32.7 & $-14.8$ & 5.04 & 3.82 & $-24.2$ \\
S4 & 3.62 & 4.26 & $+17.7$ & 53.0 & 40.7 & $-23.2$ & 5.41 & 3.25 & $-39.9$ \\
\bottomrule
\end{tabular}
\end{table}

Bajo FCFS cada courier recibe órdenes individuales (tamaño de agrupamiento igual a
1.0). Bajo RH, el tamaño promedio de agrupamiento crece monotónicamente con la
saturación, de 1.77 en S1 a 3.07 en S4 (Tabla~\ref{tab:res_bundles}), reflejando
la adaptación dinámica del algoritmo. La Figura~\ref{fig:bundle_dist} muestra que
en S1 más de la mitad de los pedidos se despachan individualmente, mientras que en
S4 la mayoría se agrupa en \textit{bundles} de 3--4. El costo por pedido converge
a $p_1=\$10$ a medida que aumenta la saturación, pues el pago variable domina
sobre el mínimo garantizado.

\begin{table}[htbp]
\centering
\small
\caption{Indicadores de agrupamiento y costo: promedio de 7 réplicas. Fuente:
elaboración propia.}
\label{tab:res_bundles}
\begin{tabular}{@{}l rr r rr r@{}}
\toprule
& \multicolumn{3}{c}{Tamaño de agrup.} & \multicolumn{3}{c}{Costo/Pedido (\$)} \\
\cmidrule(lr){2-4} \cmidrule(lr){5-7}
Esc. & FCFS & RH & $\Delta$\% & FCFS & RH & $\Delta$\% \\
\midrule
S1 & 1.00 & 1.77 & $+77.0$  & 11.40 & 11.06 & $-3.0$ \\
S2 & 1.00 & 2.28 & $+128.0$ & 10.01 & 10.07 & $+0.6$ \\
S3 & 1.00 & 2.67 & $+167.0$ & 10.00 & 10.00 & $+0.0$ \\
S4 & 1.00 & 3.07 & $+207.0$ & 10.00 & 10.00 & $+0.0$ \\
\bottomrule
\end{tabular}
\end{table}

\begin{figure}[htbp]
\centering
\includegraphics[width=0.8\textwidth]{img/fig_bundle_distribution.png}
\caption{Distribución del tamaño de agrupamiento bajo RH por escenario. Fuente:
elaboración propia.}
\label{fig:bundle_dist}
\end{figure}

\FloatBarrier

\subsection{Análisis estadístico}

La Tabla~\ref{tab:mann_whitney} presenta la prueba de Mann-Whitney~$U$ para el CtD.
Las 28 pruebas individuales (4~escenarios $\times$ 7~semillas) resultaron
significativas tras la corrección de Benjamini-Hochberg. El tamaño del efecto
disminuye de S1 a S4, lo que requiere interpretación cuidadosa: en S1--S2 el
efecto grande favorece a FCFS en la comparación ordinal; en S4, el efecto pequeño
($\bar{r}=0.144$) no refleja ausencia de diferencia, sino la alta variabilidad de
FCFS (distribución bimodal con cola de horas), pese a que RH reduce el CtD
promedio un 57.2\,\%. La prueba de Kruskal-Wallis confirma que el efecto de la
política varía significativamente entre escenarios ($H=782.78$, $p<0.001$),
justificando el análisis por regímenes de saturación.

\begin{table}[htbp]
\centering
\small
\caption{Prueba de Mann-Whitney~$U$ para el \textit{Click-to-Door}: 7 semillas
($\bar{r}\pm\sigma_r$). Fuente: elaboración propia.}
\label{tab:mann_whitney}
\begin{tabular}{@{}l rr c c l@{}}
\toprule
Esc. & $\bar{x}_{\text{FCFS}}$ & $\bar{x}_{\text{RH}}$ & $p$ & $\bar{r}\pm\sigma_r$ & Efecto \\
\midrule
S1 & 10.10 & 20.38 & $<0.001$ & $0.911 \pm 0.018$ & Grande  \\
S2 & 16.68 & 22.74 & $<0.001$ & $0.687 \pm 0.045$ & Grande  \\
S3 & 35.79 & 25.52 & $<0.001$ & $0.406 \pm 0.046$ & Mediano \\
S4 & 75.41 & 32.26 & $<0.001$ & $0.144 \pm 0.030$ & Pequeño \\
\bottomrule
\end{tabular}
\end{table}

\FloatBarrier

\subsection{Discusión}

Los resultados revelan dos regímenes con una zona de transición identificable
(Figura~\ref{fig:kpi_vs_ratio}). En el \textbf{régimen de baja saturación
(S1--S2)}, la política reactiva FCFS logra menores tiempos promedio porque
despacha cada pedido sin latencia de agrupamiento; sin embargo, ya en S2 aparece
una señal temprana de la ventaja de RH (mayor cumplimiento SLA y reducción del
\textit{CTD Overage} superior al 96\,\%). En el \textbf{régimen de saturación
media--alta (S3--S4)}, RH supera a FCFS en todas las métricas de calidad, con
reducciones del CtD del 28.7\,\% al 57.2\,\% y la eliminación total de pedidos no
entregados en S4. Las curvas de CtD se cruzan entre $r=26$ y $r=35$, marcando el
punto donde la anticipación comienza a dominar.

\begin{figure}[htbp]
\centering
\includegraphics[width=0.85\textwidth]{img/fig_kpi_vs_ratio.png}
\caption{Evolución de los KPIs principales en función de la saturación. El cruce
de las curvas de CtD entre los ratios 26 y 35 marca la transición de régimen.
Fuente: elaboración propia.}
\label{fig:kpi_vs_ratio}
\end{figure}

La Figura~\ref{fig:util_cost} muestra el compromiso entre utilización y costo por
pedido: RH logra menor utilización con \textit{throughput} equivalente o superior,
lo que confirma su mayor eficiencia kilométrica. En conjunto, estos hallazgos
superan el rango de mejora de 10--25\,\% que \citet{Reyes2018} reportan al pasar de
asignación reactiva a horizonte rodante en ciudades estadounidenses, lo que
sugiere que la topología vial de La Paz ---de menor conectividad y mayor
variabilidad de distancias--- amplifica la ventaja del agrupamiento. Este
resultado es coherente con la advertencia de \citet{Boyaci2021} y
\citet{Figliozzi2007} sobre la distorsión de las aproximaciones lineales en redes
reales.

\begin{figure}[htbp]
\centering
\includegraphics[width=0.85\textwidth]{img/fig_utilization_cost.png}
\caption{Compromiso entre utilización del repartidor y costo por pedido. Fuente:
elaboración propia.}
\label{fig:util_cost}
\end{figure}

Los resultados deben interpretarse considerando las limitaciones del modelo: la
demanda es sintética (Poisson) y puede no capturar la variabilidad de una
operación comercial; los tiempos de preparación y de servicio son homogéneos; la
red vial se asume estática (sin congestión dinámica); y tres de las diez semillas
no generaron resultados válidos por una interrupción del servicio OSRM.

% ======================================================================
\section{Conclusiones}
% ======================================================================

Se desarrolló y evaluó un entorno de simulación reproducible para el MDRP que
integra tiempos de viaje reales sobre la red vial de La Paz, B.C.S., utilizando
exclusivamente software de código abierto (OSRM, OpenStreetMap, Python y Docker).
El objetivo general se cumplió: se ejecutaron 56 corridas válidas sobre la red
real y se cuantificaron las diferencias entre las políticas FCFS y \textit{Rolling
Horizon} con pruebas estadísticas formales.

La hipótesis se verifica de forma \textbf{condicional al régimen de saturación}.
Con flota holgada (S1--S2), la política reactiva es más rápida en promedio, aunque
la anticipatoria ya mejora el cumplimiento SLA. A partir de la saturación media
(S3), RH supera a FCFS en todas las métricas de calidad de servicio, y en alta
saturación (S4) reduce el CtD promedio un 57.2\,\%, su percentil~90 un 77.1\,\% y
elimina los pedidos no entregados. El \textit{throughput} bajo RH es
consistentemente superior y la distancia por pedido se reduce hasta un 39.9\,\%.

Las contribuciones principales son: (1)~un entorno de simulación reproducible y
abierto aplicado por primera vez a una ciudad latinoamericana de escala media;
(2)~la evidencia empírica de un punto de cruce ($r\approx 26$--$35$) a partir del
cual la anticipación supera a la asignación reactiva, útil como guía práctica para
operadores; y (3)~la constatación de que los beneficios del horizonte rodante no
solo se mantienen sino que se amplifican sobre una topología vial real y hostil,
superando las mejoras reportadas en la literatura para redes de cuadrícula. Al
prescindir de APIs comerciales, el enfoque ofrece una alternativa de soberanía
tecnológica que reduce las barreras de entrada para las PyMES locales
\citep{MunozVillamizar2021}.

Como líneas de trabajo futuro se plantea: diagnosticar la condición de borde que
invalidó tres semillas; sustituir la demanda sintética por datos operacionales
reales; incorporar perfiles de congestión dinámica en OSRM; y localizar con mayor
precisión el punto de cruce explorando valores intermedios de $r$.

% ======================================================================
\bibliographystyle{apacite}
\bibliography{ref/biblio}

\end{document}
